\subsection{The normalized level-zero word energy}
\label{sec:archimedean-spectral-word-energy}

We use the original positive Jacquet pairing \(\beta\) and the
exact level-zero word identity, with all Haar and compact-type
normalizations unchanged. The inputs from the construction record
are the signed standard and dual-standard Pieri decompositions
and the mixed Jacquet orthogonality formula. Their proofs are in
\texttt{global-norm-frontier-pieri-family.tex} and
\texttt{global-norm-frontier-jacquet-gram.tex}.\cite{current-construction-record}
We retain their smooth-decay, original period-normalization, and
shift-reserve hypotheses. Orthogonality will eliminate dependence
on the number or the norms of the tensor projections.

Let \(s=r+1\), \(n=2r+k+1\), and let \(F\) be a finite smooth
coefficient of the original discrete series of real
Harish--Chandra parameter \(\lambda'=(\lambda'_j)_{j=0}^{n-1}\).
Let \(A(g)\) be a finite operator-valued polynomial of degree at
most \(d_A\) in the entries of \(g\) and their conjugates.
For an integer \(p\ge0\) put
\[
                  H(g)=G(g)^pF(g)A(g),\qquad
                         d=2sp+d_A .                  \tag{SE1}
\]
The Gram polynomial has degree \(2s\), so the polynomial
multiplier has degree at most \(d\).

\begin{theorem}[Same-form spectral bound]
\label{thm:archimedean-spectral-word-energy}
Assume that the original signed standard and dual-standard
families and mixed Jacquet identity are available throughout
the shift ball \(\|\eta\|_1\le d\) about \(\lambda'\), that
the parameters stay in the same regular chamber, and that the
original smooth-decay reserve supplies absolute convergence of
the pairings below. Let \(\lambda\) be the base parameter and
\(c\ge1\) satisfy
\[
                  |\lambda'_j-\lambda_j|+d\le c-1
                      \quad(0\le j<n).
                                                        \tag{SE2}
\]
Define
\[
 \nu(t)=\prod_{j=0}^{n-1}(|t+\lambda_j|+c),\qquad
                 S_w=\prod_{(\epsilon,t)\in w}\nu(t).
                                                        \tag{SE3}
\]
Then, for every admissible original word \(w\),
\[
                      E_{0,w}(H)\le\beta(H,H)=jN^2(H).
                                                        \tag{SE4}
\]
There is no coefficient count or compact-type dimension in
(SE4). The conclusion includes \(jN^2(H)=0\).
\end{theorem}

\begin{proof}
Each entry of \(g\) is a coefficient of the standard
representation. On the unitary group, each conjugate entry is
a linear combination of coefficients of the dual standard
representation, since \(g^{-1}=Jg^*J\) for the original
Hermitian matrix \(J\). Hence every multiplier monomial in
(SE1) is a coefficient of a tensor product of at most \(d\)
standard or dual-standard factors. Apply the signed one-step
Pieri decomposition successively in the original family.
Each step changes one parameter coordinate by \(+1\) or \(-1\);
all intermediate shifts have \(\ell^1\)-norm at most \(d\)
and are allowed by the stated reserve.

This expresses \(H\) as a finite sum of smooth coefficients
with parameters \(\lambda'+\eta\), \(\|\eta\|_1\le d\).
Combine all summands with the same final parameter:
\[
                          H=\sum_{\eta\in\mathcal S}H_\eta.
                                                        \tag{SE5}
\]
Distinct parameters in this fixed chamber have distinct central
characters. The mixed Jacquet identity gives
\[
 \begin{split}
 \beta(H_\eta,H_\zeta)&=0\quad(\eta\ne\zeta),\\
 \beta(H,H)&=\sum_{\eta\in\mathcal S}\beta(H_\eta,H_\eta),
 \qquad \beta(H_\eta,H_\eta)\ge0.
 \end{split}                                           \tag{SE6}
\]
The signed tensor form obtains the coefficients; the form in
(SE6) is the original positive radial pairing. No bound on the
indefinite tensor projections is needed.

The full central Capelli operator \(\mathscr C(t)\) acts on
\(H_\eta\) by the real scalar
\[
                   p_\eta(t)=\prod_j(t+\lambda'_j+\eta_j).
\]
The exact empty-head identity is
\[
 E_{0,w}(H)=S_w^{-1}
      \operatorname{Re}\beta(\mathscr C_w H,H),
                                                        \tag{SE7}
\]
where \(\mathscr C_w\) is the ordered product of the full
Capelli operators at the original numerical letter parameters,
with their prescribed real signs. The signs have modulus one;
no letter shift or word normalizer is omitted. Its eigenvalue
on \(H_\eta\) is a real number \(a_{w,\eta}\) with
\[
 |a_{w,\eta}|=\prod_{(\epsilon,t)\in w}|p_\eta(t)|
                                      \le S_w .         \tag{SE8}
\]
Indeed, (SE2) and \(|\eta_j|\le d\) give
\[
 |t+\lambda'_j+\eta_j|
       \le |t+\lambda_j|+|\lambda'_j-\lambda_j|+d
       \le |t+\lambda_j|+c.
\]
It follows from (SE6)--(SE8) that
\[
 \operatorname{Re}\beta(\mathscr C_wH,H)
      =\sum_\eta a_{w,\eta}\beta(H_\eta,H_\eta)
      \le S_w\sum_\eta\beta(H_\eta,H_\eta).
\]
Divide by \(S_w>0\) and use the unchanged empty-word radial
identity \(\beta(H,H)=jN^2(H)\). If this norm vanishes,
every nonnegative term in (SE6) vanishes, so the argument
also proves the zero-norm case.
\end{proof}

This theorem bounds the uncut coefficient. It does not itself
justify differentiating a Gram cutoff and then removing it.
In particular it cannot be applied directly to
\(\chi(G/\varepsilon)H\), whose nonpolynomial cutoff does
not have the finite spectral decomposition (SE5).
